% GENERATED FILE. Edit canonical lessons, then run npm run generate:books.
% canonical-source-hash: fnv1a64:ca7234e13aba8055
% canonical-lessons: JA-C75-maitoshi, JA-C75-kotoshi, JA-C75-konshuu, JA-C75-ototoshi, JA-C75-asanebou

\chapter{Every year, This year, This week, Year before last, Oversleeping}
\label{ch:ja-every-year-this-year-this-week-year-before-last-oversleeping}
\hlchaptermodality{75}

\begin{chapteropening}
\textbf{By the end of this chapter:} \emph{I can say every year, this year, this week, the year before last and oversleeping.}

Complete the last lesson of chapter 75: i can say every year, this year, this week, the year before last and oversleeping.
\end{chapteropening}

\section[\texorpdfstring{\ja{まいとし} (maitoshi)}{まいとし (maitoshi)}]{\ja{まいとし} (maitoshi) — every year}

\label{lesson:JA-C75-maitoshi}



\begin{quote}
Before the new one: say the Japanese for on the way, then the Japanese for half past.
\end{quote}

\subsection*{You'll want to know: \ja{まいとし}}
\textbf{\ja{まいとし}} — \emph{maitoshi} — \textquotedblleft{}every year\textquotedblright{}.

\begin{grammarlens}[title={What you've built}]
One more word for this chapter.
\end{grammarlens}

\subsection*{Your turn}
\begin{itemize}
\raggedright
  \item \emph{Say it:} \emph{maitoshi}
  \item \emph{Say it:} \emph{maitoshi}, once more
  \item \emph{Say it:} say \emph{maitoshi}
  \item \emph{Recall:} say the Japanese for on the way, then the Japanese for half past, then say \emph{maitoshi} again
\end{itemize}

\begin{tcolorbox}[breakable,colback=teal!4,colframe=teal!35!black,title={Before you move on}]
What does \ja{まいとし} mean? (\textquotedblleft{}Every year\textquotedblright{}.) Say it once more.
\end{tcolorbox}
\section[\texorpdfstring{\ja{ことし} (kotoshi)}{ことし (kotoshi)}]{\ja{ことし} (kotoshi) — this year}

\label{lesson:JA-C75-kotoshi}



\begin{quote}
Before the new one: say the Japanese for half past, then the Japanese for every year.
\end{quote}

\subsection*{You'll want to know: \ja{ことし}}
\textbf{\ja{ことし}} — \emph{kotoshi} — \textquotedblleft{}this year\textquotedblright{}.

\begin{grammarlens}[title={What you've built}]
One more word for this chapter.
\end{grammarlens}

\subsection*{Your turn}
\begin{itemize}
\raggedright
  \item \emph{Say it:} \emph{kotoshi}
  \item \emph{Say it:} \emph{kotoshi}, once more
  \item \emph{Say it:} say \emph{kotoshi}
  \item \emph{Recall:} say the Japanese for half past, then the Japanese for every year, then say \emph{kotoshi} again
\end{itemize}

\begin{tcolorbox}[breakable,colback=teal!4,colframe=teal!35!black,title={Before you move on}]
What does \ja{ことし} mean? (\textquotedblleft{}This year\textquotedblright{}.) Say it once more.
\end{tcolorbox}
\section[\texorpdfstring{\ja{こんしゅう} (konshuu)}{こんしゅう (konshuu)}]{\ja{こんしゅう} (konshuu) — this week}

\label{lesson:JA-C75-konshuu}



\begin{quote}
Before the new one: say the Japanese for every year, then the Japanese for this year.
\end{quote}

\subsection*{You'll want to know: \ja{こんしゅう}}
\textbf{\ja{こんしゅう}} — \emph{konshuu} — \textquotedblleft{}this week\textquotedblright{}.

\begin{grammarlens}[title={What you've built}]
One more word for this chapter.
\end{grammarlens}

\subsection*{Your turn}
\begin{itemize}
\raggedright
  \item \emph{Say it:} \emph{konshuu}
  \item \emph{Say it:} \emph{konshuu}, once more
  \item \emph{Say it:} say \emph{konshuu}
  \item \emph{Recall:} say the Japanese for every year, then the Japanese for this year, then say \emph{konshuu} again
\end{itemize}

\begin{tcolorbox}[breakable,colback=teal!4,colframe=teal!35!black,title={Before you move on}]
What does \ja{こんしゅう} mean? (\textquotedblleft{}This week\textquotedblright{}.) Say it once more.
\end{tcolorbox}
\section[\texorpdfstring{\ja{おととし} (ototoshi)}{おととし (ototoshi)}]{\ja{おととし} (ototoshi) — the year before last}

\label{lesson:JA-C75-ototoshi}



\begin{quote}
Before the new one: say the Japanese for this year, then the Japanese for this week.
\end{quote}

\subsection*{You'll want to know: \ja{おととし}}
\textbf{\ja{おととし}} — \emph{ototoshi} — \textquotedblleft{}the year before last\textquotedblright{}.

\begin{grammarlens}[title={What you've built}]
One more word for this chapter.
\end{grammarlens}

\subsection*{Your turn}
\begin{itemize}
\raggedright
  \item \emph{Say it:} \emph{ototoshi}
  \item \emph{Say it:} \emph{ototoshi}, once more
  \item \emph{Say it:} say \emph{ototoshi}
  \item \emph{Recall:} say the Japanese for this year, then the Japanese for this week, then say \emph{ototoshi} again
\end{itemize}

\begin{tcolorbox}[breakable,colback=teal!4,colframe=teal!35!black,title={Before you move on}]
What does \ja{おととし} mean? (\textquotedblleft{}The year before last\textquotedblright{}.) Say it once more.
\end{tcolorbox}
\section[\texorpdfstring{\ja{あさねぼう} (asanebou)}{あさねぼう (asanebou)}]{\ja{あさねぼう} (asanebou) — oversleeping}

\label{lesson:JA-C75-asanebou}



\begin{quote}
Before the new one: say the Japanese for this week, then the Japanese for the year before last.
\end{quote}

\subsection*{You'll want to know: \ja{あさねぼう}}
\textbf{\ja{あさねぼう}} — \emph{asanebou} — \textquotedblleft{}oversleeping\textquotedblright{}.

\begin{grammarlens}[title={What you've built}]
One more word for this chapter.
\end{grammarlens}

\subsection*{Your turn}
\begin{itemize}
\raggedright
  \item \emph{Say it:} \emph{asanebou}
  \item \emph{Say it:} \emph{asanebou}, once more
  \item \emph{Say it:} say \emph{asanebou}
  \item \emph{Recall:} say the Japanese for this week, then the Japanese for the year before last, then say \emph{asanebou} again
\end{itemize}

\begin{tcolorbox}[breakable,colback=teal!4,colframe=teal!35!black,title={Before you move on}]
What does \ja{あさねぼう} mean? (\textquotedblleft{}Oversleeping\textquotedblright{}.) Say it once more.
\end{tcolorbox}
